\documentclass[10pt]{scrartcl}
\usepackage[mdthm,colorsec,noauthor,noasy]{evan}
\usepackage{handout}
\DeclareMathOperator{\im}{range}
\title{MATH 110 --- Week 6}
\subtitle{Homework 6 --- Solutions}
\author{}\date{}\ihead{}
\begin{document}\maketitle
\begin{prob}\textbf{HW6 P1.}\par\idea{The equation $P^2=P$ controls eigenvalues and supplies a decomposition for every vector.}\begin{steps}\ii Assume the intended linearity hypothesis. If $Pv=\lambda v$ with $v\ne0$, then $P^2v=\lambda^2v$, while $P^2v=Pv=\lambda v$. Hence $\lambda(\lambda-1)=0$, so $\lambda\in\{0,1\}$. It is not asserted that both occur.\ii Every $v=(v-Pv)+Pv$, with $P(v-Pv)=Pv-P^2v=0$ and $Pv\in\im P$. Thus $V=\ker P+\im P$.\ii If $w\in\ker P\cap\im P$, write $w=Pu$. Then $w=Pu=P^2u=Pw=0$, so the sum is direct.\end{steps}\end{prob}
\begin{prob}\textbf{HW6 P2(a).}\par\idea{Differentiation lowers the degree of every nonconstant polynomial over $\mathbb R$ or $\mathbb C$.}\begin{steps}\ii If $Dp=\lambda p$ and $p\ne0$ with $\lambda\ne0$, a nonconstant $p$ would give degree $\deg p-1$ on the left and $\deg p$ on the right, a contradiction. A nonzero constant gives zero on the left and nonzero on the right, also impossible.\ii The eigenvalue $0$ occurs because $D1=0$. Writing $p=\sum_{j=0}^n a_jx^j$ gives $Dp=\sum_{j=1}^n ja_jx^{j-1}$. A zero derivative forces each $ja_j=0$, hence $a_j=0$ for $j\ge1$ in either field. Thus the kernel consists of constants.\ii Thus the only eigenvalue is 0 and $E(0,D)=\operatorname{span}(1)$; for $\lambda\ne0$, $E(\lambda,D)=\{0\}$. This also covers $n=0$.\end{steps}\end{prob}
\begin{prob}\textbf{HW6 P2(b).}\par\idea{Smooth functions include nonpolynomial exponential solutions of the derivative equation.}\begin{steps}\ii We cite elementary calculus: exponential differentiation, the product rule, and the fact that a function with derivative zero on $\mathbb R$ is constant. An eigenfunction is a nonzero function satisfying $Df=\lambda f$. On smooth real-valued functions, for every real $\lambda$ the function $e^{\lambda x}$ is a nonzero eigenfunction because its derivative is $\lambda e^{\lambda x}$.\ii If $f^{\prime}=\lambda f$, then $(e^{-\lambda x}f(x))^{\prime}=0$, so $f=ce^{\lambda x}$. Thus $E(\lambda,D)=\operatorname{span}(e^{\lambda x})$.\ii A complex-valued function is smooth when its real and imaginary parts are smooth; differentiation acts on those parts separately. For these functions, the same argument applies for every complex $\lambda=a+ib$, using $e^{\lambda x}=e^{ax}(\cos bx+i\sin bx)$ and the elementary trigonometric derivative rules. Hence all scalars of the selected field occur, unlike the polynomial case.\end{steps}\end{prob}
\begin{prob}\textbf{HW6 P3.}\par\idea{The invertible map $S$ sends the equation for $R$ to the equation for $T$.}\begin{steps}\ii $Rv=\lambda v$ iff $S^{-1}TSv=\lambda v$, iff $T(Sv)=\lambda Sv$ after multiplication by $S$.\ii Thus $v\in E(\lambda,R)$ iff $Sv\in E(\lambda,T)$, proving $E(\lambda,R)=S^{-1}E(\lambda,T)$.\ii Since $S$ is invertible, it sends nonzero vectors to nonzero vectors and back. Therefore one eigenspace is nonzero iff the other is, giving equal eigenvalues.\end{steps}\end{prob}
\begin{prob}\textbf{HW6 P4.}\par\idea{For two-by-two matrices, noninvertibility is exactly a zero elementary determinant.}\begin{steps}\ii $tI-A=\left(\begin{smallmatrix}t-a&-b\\-c&t-d\end{smallmatrix}\right)$, so its determinant is $(t-a)(t-d)-bc=t^2-(a+d)t+(ad-bc)$.\ii $\lambda$ is an eigenvalue iff $(\lambda I-A)v=0$ for nonzero $v$, iff $\lambda I-A$ is not injective. Finite dimension makes this equivalent to noninvertibility.\ii Week 4's two-by-two criterion makes that equivalent to $\det(\lambda I-A)=0$.\end{steps}\end{prob}
\begin{prob}\textbf{HW6 P5: $A_1$.}\par\idea{First specify the field, then solve the eigenvalue equation.}\begin{steps}\ii $\chi_{A_1}(t)=t^2+1$, so over $\mathbb R$ there are no eigenvalues and no eigenvector basis.\ii Over $\mathbb C$ the eigenvalues are $i,-i$. The equations $-y=\lambda x,x=\lambda y$ give $E(i,A_1)=\mathbb C(1,-i)$ and $E(-i,A_1)=\mathbb C(1,i)$.\ii The two eigenvalues are distinct, so their displayed eigenvectors are independent by Proposition 54 and form a basis of $\mathbb C^2$.\end{steps}\end{prob}
\begin{prob}\textbf{HW6 P5: $A_2,A_3$.}\par\idea{A root gives a candidate eigenvalue; the kernel dimension determines how many independent eigenvectors it supplies.}\begin{steps}\ii For $A_2$, $\chi(t)=(t-1)^2$ and $(A_2-I)(x,y)=(y,0)$. Thus the sole eigenvalue is 1, with $E(1,A_2)=\mathbb Re_1$. A one-dimensional eigenspace cannot supply a basis of $\mathbb R^2$.\ii For $A_3$, $\chi(t)=t^2-t-1$, with $\lambda_\pm=(1\pm\sqrt5)/2$. Solving gives $E(\lambda_\pm,A_3)=\mathbb R(1,\lambda_\pm)$.\ii These eigenvalues are distinct, so the two displayed eigenvectors form a basis of $\mathbb R^2$ by Proposition 54.\end{steps}\end{prob}
\begin{prob}\textbf{HW6 P6 (optional).}\par\idea{The first dependent iterate makes a finite initial span invariant.}\begin{steps}\ii A vector of $C(v,T)$ is a finite sum $\sum_{j=0}^r a_jT^jv$. Applying $T$ gives $\sum_{j=0}^r a_jT^{j+1}v$, again a finite combination of the iterates. Thus the cyclic subspace is invariant.\ii Let $n=\dim V$. Since $v\ne0$ and the first $n+1$ iterates are dependent, there is a first index $m$ with $1\le m\le n$ such that $T^mv\in\operatorname{span}(v,\ldots,T^{m-1}v)$.\ii The span $U$ of those first $m$ iterates is invariant: $T$ sends each generator except the last to the next, and sends the last into $U$ by the first-dependence relation.\ii Induction gives every later iterate in $U$, so $C(v,T)=U$, as required.\end{steps}\end{prob}
\begin{prob}\textbf{HW6 P7 (optional).}\par\idea{Stationarity forces equal coordinates, and total probability fixes their value.}\begin{steps}\ii Write $p=(a,b,c)^{\mathsf T}$. The first equation $a=(a+b)/2$ gives $a=b$; the third $c=(b+c)/2$ gives $c=b$.\ii The second equation $b=(a+c)/2$ then holds. Since $a+b+c=1$, the probability vector is $p=(1/3,1/3,1/3)^{\mathsf T}$.\ii Thus if the chain starts uniformly over the three states, its distribution remains uniform after each transition.\end{steps}\end{prob}
\begin{prob}\textbf{Exam practice, Problem 7.}\par\idea{Invariance passes to annihilators because precomposition evaluates only on images that stay in the subspace.}\begin{steps}\ii If $v\in\ker T$, then $Tv=0\in\ker T$; if $w=Tv\in\im T$, then $Tw=T(Tv)\in\im T$. Thus both subspaces are invariant.\ii More generally, let $U$ be invariant and $\varphi\in U^0$. For each $u\in U$, $Tu\in U$, so $(T^{\prime}\varphi)(u)=\varphi(Tu)=0$.\ii Therefore $T^{\prime}\varphi\in U^0$, proving annihilator invariance. Apply this to $U=\ker T$ and $U=\im T$.\end{steps}\end{prob}\end{document}
